\documentclass[11pt,a4paper]{article}
\usepackage{fontspec}
\setmainfont{Times New Roman}
\usepackage[top=1.5cm, bottom=1.5cm, left=1.8cm, right=1.8cm]{geometry}
\usepackage{amsmath,amssymb}
\usepackage{fvextra, xcolor}

\pagestyle{empty}

\DefineVerbatimEnvironment{htmlcode}{Verbatim}{
  breaklines=true,
  breaksymbolleft={},
  breaksymbolright={},
  frame=none,
  fontsize=\fontsize{8pt}{10pt}\selectfont,
  samepage=true
}

\newcommand{\fourchoices}[4]{%
  \par\noindent\hspace*{0.3cm}%
  \makebox[0.245\dimexpr\textwidth-0.3cm\relax][l]{\textbf{[A. ]} #1}%
  \makebox[0.245\dimexpr\textwidth-0.3cm\relax][l]{\textbf{B. } #2}%
  \makebox[0.245\dimexpr\textwidth-0.3cm\relax][l]{\textbf{C. } #3}%
  \makebox[0.245\dimexpr\textwidth-0.3cm\relax][l]{\textbf{D. } #4}\par\vspace{5pt}%
}

\newcommand{\twochoices}[4]{%
  \par\noindent\hspace*{0.3cm}%
  \makebox[0.49\dimexpr\textwidth-0.3cm\relax][l]{\textbf{[A. ]} #1}%
  \makebox[0.49\dimexpr\textwidth-0.4cm\relax][l]{\textbf{B. } #2}\par\vspace{2.5pt}%
  \noindent\hspace*{0.3cm}%
  \makebox[0.49\dimexpr\textwidth-0.3cm\relax][l]{\textbf{C. } #3}%
  \makebox[0.49\dimexpr\textwidth-0.4cm\relax][l]{\textbf{D. } #4}\par\vspace{5pt}%
}

\newcommand{\onechoice}[4]{%
  \par
  {\hangindent=1.1cm\hangafter=1\noindent\hspace*{0.3cm}\textbf{[A. ]} #1\par\vspace{2pt}}%
  {\hangindent=1.1cm\hangafter=1\noindent\hspace*{0.3cm}\textbf{B. } #2\par\vspace{2pt}}%
  {\hangindent=1.1cm\hangafter=1\noindent\hspace*{0.3cm}\textbf{C. } #3\par\vspace{2pt}}%
  {\hangindent=1.1cm\hangafter=1\noindent\hspace*{0.3cm}\textbf{D. } #4\par\vspace{5pt}}%
}

\newlength{\choiceawidth}\newlength{\choicebwidth}\newlength{\choicecwidth}\newlength{\choicedwidth}\newlength{\maxchoicewidth}
\newcommand{\choices}[4]{%
  \settowidth{\choiceawidth}{\textbf{[A. ]} #1}%
  \settowidth{\choicebwidth}{\textbf{B. } #2}%
  \settowidth{\choicecwidth}{\textbf{C. } #3}%
  \settowidth{\choicedwidth}{\textbf{D. } #4}%
  \setlength{\maxchoicewidth}{\choiceawidth}%
  \ifdim\choicebwidth>\maxchoicewidth \setlength{\maxchoicewidth}{\choicebwidth}\fi
  \ifdim\choicecwidth>\maxchoicewidth \setlength{\maxchoicewidth}{\choicecwidth}\fi
  \ifdim\choicedwidth>\maxchoicewidth \setlength{\maxchoicewidth}{\choicedwidth}\fi
  \ifdim\maxchoicewidth<0.23\dimexpr\textwidth-0.3cm\relax
    \fourchoices{#1}{#2}{#3}{#4}%
  \else
    \ifdim\maxchoicewidth<0.47\dimexpr\textwidth-0.3cm\relax
      \twochoices{#1}{#2}{#3}{#4}%
    \else
      \onechoice{#1}{#2}{#3}{#4}%
    \fi
  \fi
}

\begin{document}

\noindent\textbf{Câu 1.} Quan sát đoạn mã lệnh sau và cho biết kết quả hiển thị trên màn hình:
\begin{htmlcode}
<div style="background-color: #0f172a; color: #ffffff; padding: 10px 16px; margin: 6px 0; border-radius: 6px; border: 1.5px solid #3b82f6; font-family: monospace; font-size: 14px;"><span style="color: #38bdf8; font-weight: bold;">print</span><span style="color: #ffffff; font-weight: bold;">(</span><span style="color: #4ade80; font-weight: bold;">15</span> <span style="color: #f472b6; font-weight: bold;">/</span> <span style="color: #4ade80; font-weight: bold;">3</span><span style="color: #ffffff; font-weight: bold;">)</span></div>
\end{htmlcode}
\choices{\texttt{5.0}}{\texttt{5}}{\texttt{3}}{Báo lỗi cú pháp}

\noindent\textbf{Câu 2.} Cho đoạn mã lệnh sau. Kết quả xuất ra màn hình là gì?
\begin{htmlcode}
<div style="background-color: #0f172a; color: #ffffff; padding: 10px 16px; margin: 6px 0; border-radius: 6px; border: 1.5px solid #3b82f6; font-family: monospace; font-size: 14px;"><span style="color: #38bdf8; font-weight: bold;">print</span><span style="color: #ffffff; font-weight: bold;">(</span><span style="color: #fde047; font-weight: bold;">"3 + 7 ="</span>, <span style="color: #4ade80; font-weight: bold;">3</span> <span style="color: #f472b6; font-weight: bold;">+</span> <span style="color: #4ade80; font-weight: bold;">7</span><span style="color: #ffffff; font-weight: bold;">)</span></div>
\end{htmlcode}
\choices{\texttt{3 + 7 = 10}}{\texttt{10 = 10}}{\texttt{3 + 7 = 3 + 7}}{\texttt{3 + 7 = 37}}

\noindent\textbf{Câu 3.} Cho đoạn mã lệnh sau:
\begin{htmlcode}
<div style="background-color: #0f172a; color: #ffffff; padding: 10px 16px; margin: 6px 0; border-radius: 6px; border: 1.5px solid #3b82f6; font-family: monospace; font-size: 14px;"><span style="color: #38bdf8; font-weight: bold;">print</span><span style="color: #ffffff; font-weight: bold;">(</span><span style="color: #fde047; font-weight: bold;">"Điểm Tin:"</span>, <span style="color: #4ade80; font-weight: bold;">10</span>, <span style="color: #fde047; font-weight: bold;">"học sinh:"</span>, <span style="color: #4ade80; font-weight: bold;">40</span><span style="color: #ffffff; font-weight: bold;">)</span></div>
\end{htmlcode}
Màn hình sẽ hiển thị kết quả nào dưới đây?
\choices{\texttt{Điểm Tin: 10 học sinh: 40}}{\texttt{Điểm Tin:10học sinh:40}}{\texttt{Điểm Tin:, 10, học sinh:, 40}}{Báo lỗi vì không thể in chữ kèm số}

\noindent\textbf{Câu 4.} Cho đoạn chương trình sau. Kết quả hiển thị trên màn hình là bao nhiêu?
\begin{htmlcode}
<div style="background-color: #0f172a; color: #ffffff; padding: 10px 16px; margin: 6px 0; border-radius: 6px; border: 1.5px solid #3b82f6; font-family: monospace; font-size: 14px;"><span style="color: #38bdf8; font-weight: bold;">print</span><span style="color: #ffffff; font-weight: bold;">(</span><span style="color: #4ade80; font-weight: bold;">20</span> <span style="color: #f472b6; font-weight: bold;">-</span> <span style="color: #4ade80; font-weight: bold;">4</span> <span style="color: #f472b6; font-weight: bold;">*</span> <span style="color: #4ade80; font-weight: bold;">3</span><span style="color: #ffffff; font-weight: bold;">)</span></div>
\end{htmlcode}
\choices{\texttt{8}}{\texttt{48}}{\texttt{8.0}}{\texttt{16}}

\noindent\textbf{Câu 5.} Quan sát đoạn mã lệnh sau và cho biết kết quả xuất ra màn hình:
\begin{htmlcode}
<div style="background-color: #0f172a; color: #ffffff; padding: 10px 16px; margin: 6px 0; border-radius: 6px; border: 1.5px solid #3b82f6; font-family: monospace; font-size: 14px;"><span style="color: #38bdf8; font-weight: bold;">print</span><span style="color: #ffffff; font-weight: bold;">(</span><span style="color: #ffffff; font-weight: bold;">(</span><span style="color: #4ade80; font-weight: bold;">10</span> <span style="color: #f472b6; font-weight: bold;">+</span> <span style="color: #4ade80; font-weight: bold;">6</span><span style="color: #ffffff; font-weight: bold;">)</span> <span style="color: #f472b6; font-weight: bold;">/</span> <span style="color: #4ade80; font-weight: bold;">2</span><span style="color: #ffffff; font-weight: bold;">)</span></div>
\end{htmlcode}
\choices{\texttt{8.0}}{\texttt{8}}{\texttt{13.0}}{\texttt{13}}

\noindent\textbf{Câu 6.} Cho đoạn lệnh thực hiện phép tính giữa số nguyên và số thực:
\begin{htmlcode}
<div style="background-color: #0f172a; color: #ffffff; padding: 10px 16px; margin: 6px 0; border-radius: 6px; border: 1.5px solid #3b82f6; font-family: monospace; font-size: 14px;"><span style="color: #38bdf8; font-weight: bold;">print</span><span style="color: #ffffff; font-weight: bold;">(</span><span style="color: #4ade80; font-weight: bold;">12</span> <span style="color: #f472b6; font-weight: bold;">+</span> <span style="color: #4ade80; font-weight: bold;">1.5</span><span style="color: #ffffff; font-weight: bold;">)</span></div>
\end{htmlcode}
Kết quả hiển thị trên màn hình là:
\choices{\texttt{13.5}}{\texttt{13}}{\texttt{121.5}}{Báo lỗi vì khác kiểu dữ liệu}

\noindent\textbf{Câu 7.} Cho đoạn mã lệnh sau:
\begin{htmlcode}
<div style="background-color: #0f172a; color: #ffffff; padding: 10px 16px; margin: 6px 0; border-radius: 6px; border: 1.5px solid #3b82f6; font-family: monospace; font-size: 14px;"><span style="color: #38bdf8; font-weight: bold;">print</span><span style="color: #ffffff; font-weight: bold;">(</span><span style="color: #4ade80; font-weight: bold;">10</span> <span style="color: #f472b6; font-weight: bold;">+</span> <span style="color: #4ade80; font-weight: bold;">7</span> <span style="color: #f472b6; font-weight: bold;">/</span> <span style="color: #4ade80; font-weight: bold;">2</span><span style="color: #ffffff; font-weight: bold;">)</span></div>
\end{htmlcode}
Kết quả hiển thị trên màn hình là bao nhiêu?
\choices{\texttt{13.5}}{\texttt{13}}{\texttt{8.5}}{\texttt{14.0}}

\noindent\textbf{Câu 8.} Quan sát đoạn chương trình sau:
\begin{htmlcode}
<div style="background-color: #0f172a; color: #ffffff; padding: 10px 16px; margin: 6px 0; border-radius: 6px; border: 1.5px solid #3b82f6; font-family: monospace; font-size: 14px;"><span style="color: #94a3b8; font-style: italic;"># print("Chào buổi sáng")</span><br /><span style="color: #38bdf8; font-weight: bold;">print</span><span style="color: #ffffff; font-weight: bold;">(</span><span style="color: #fde047; font-weight: bold;">"Chào buổi chiều"</span><span style="color: #ffffff; font-weight: bold;">)</span></div>
\end{htmlcode}
Khi chạy chương trình trên, màn hình sẽ hiển thị như thế nào?
\choices{Chỉ in ra màn hình: \texttt{Chào buổi chiều}}{In ra cả 2 dòng: \texttt{Chào buổi sáng} và \texttt{Chào buổi chiều}}{Chỉ in ra màn hình: \texttt{Chào buổi sáng}}{Báo lỗi cú pháp do có dấu \texttt{\#}}

\noindent\textbf{Câu 9.} Cho đoạn mã nguồn gồm 2 câu lệnh liên tiếp:
\begin{htmlcode}
<div style="background-color: #0f172a; color: #ffffff; padding: 10px 16px; margin: 6px 0; border-radius: 6px; border: 1.5px solid #3b82f6; font-family: monospace; font-size: 14px;"><span style="color: #38bdf8; font-weight: bold;">print</span><span style="color: #ffffff; font-weight: bold;">(</span><span style="color: #fde047; font-weight: bold;">"Lớp 10A1"</span><span style="color: #ffffff; font-weight: bold;">)</span><br /><span style="color: #38bdf8; font-weight: bold;">print</span><span style="color: #ffffff; font-weight: bold;">(</span><span style="color: #fde047; font-weight: bold;">"Trường Á Châu"</span><span style="color: #ffffff; font-weight: bold;">)</span></div>
\end{htmlcode}
Kết quả hiển thị sẽ như thế nào?
\choices{In \texttt{Lớp 10A1} ở dòng 1 và \texttt{Trường Á Châu} ở dòng 2}{In \texttt{Lớp 10A1 Trường Á Châu} trên cùng một dòng}{In \texttt{Lớp 10A1Trường Á Châu} dính liền nhau}{Báo lỗi vì không được viết 2 lệnh \texttt{print} liên tiếp}

\noindent\textbf{Câu 10.} Khi gõ và thực thi câu lệnh sau trong Python:
\begin{htmlcode}
<div style="background-color: #0f172a; color: #ffffff; padding: 10px 16px; margin: 6px 0; border-radius: 6px; border: 1.5px solid #3b82f6; font-family: monospace; font-size: 14px;"><span style="color: #ef4444; font-weight: bold;">Print</span><span style="color: #ffffff; font-weight: bold;">(</span><span style="color: #fde047; font-weight: bold;">"Xin chào các bạn"</span><span style="color: #ffffff; font-weight: bold;">)</span></div>
\end{htmlcode}
Chương trình sẽ đưa ra phản hồi nào?
\choices{Báo lỗi \texttt{NameError} vì Python phân biệt chữ hoa, chữ thường}{In ra bình thường dòng chữ: \texttt{Xin chào các bạn}}{Báo lỗi \texttt{SyntaxError} do thiếu dấu chấm phẩy}{Tự động chuyển thành chữ thường và chạy bình thường}

\noindent\textbf{Câu 11.} Khi chạy câu lệnh sau trên trình soạn thảo Python:
\begin{htmlcode}
<div style="background-color: #0f172a; color: #ffffff; padding: 10px 16px; margin: 6px 0; border-radius: 6px; border: 1.5px solid #3b82f6; font-family: monospace; font-size: 14px;"><span style="color: #38bdf8; font-weight: bold;">print</span><span style="color: #ffffff; font-weight: bold;">(</span><span style="color: #fde047; font-weight: bold;">"Học lập trình Python</span><span style="color: #ffffff; font-weight: bold;">)</span></div>
\end{htmlcode}
Phần mềm Python sẽ thông báo loại lỗi nào?
\choices{\texttt{SyntaxError: unterminated string literal}}{\texttt{ZeroDivisionError}}{\texttt{TypeError}}{\texttt{IndexError}}

\noindent\textbf{Câu 12.} Cho đoạn mã lệnh sau:
\begin{htmlcode}
<div style="background-color: #0f172a; color: #ffffff; padding: 10px 16px; margin: 6px 0; border-radius: 6px; border: 1.5px solid #3b82f6; font-family: monospace; font-size: 14px;"><span style="color: #38bdf8; font-weight: bold;">print</span><span style="color: #ffffff; font-weight: bold;">(</span><span style="color: #fde047; font-weight: bold;">"Kết quả là:"</span> <span style="color: #4ade80; font-weight: bold;">5</span> <span style="color: #f472b6; font-weight: bold;">*</span> <span style="color: #4ade80; font-weight: bold;">10</span><span style="color: #ffffff; font-weight: bold;">)</span></div>
\end{htmlcode}
Vì sao đoạn lệnh trên bị báo lỗi cú pháp (\texttt{SyntaxError})?
\choices{Vì thiếu dấu phẩy \texttt{,} ngăn cách giữa xâu ký tự và biểu thức}{Vì hàm \texttt{print} không tính được phép nhân \texttt{5 * 10}}{Vì không được đặt chữ và số trong cùng hàm \texttt{print}}{Vì biểu thức \texttt{5 * 10} bắt buộc phải đặt trong dấu ngoặc kép}

\noindent\textbf{Câu 13.} Cho câu lệnh sau. Màn hình sẽ xuất hiện nội dung gì?
\begin{htmlcode}
<div style="background-color: #0f172a; color: #ffffff; padding: 10px 16px; margin: 6px 0; border-radius: 6px; border: 1.5px solid #3b82f6; font-family: monospace; font-size: 14px;"><span style="color: #38bdf8; font-weight: bold;">print</span><span style="color: #ffffff; font-weight: bold;">(</span><span style="color: #fde047; font-weight: bold;">"10 + 20"</span><span style="color: #ffffff; font-weight: bold;">)</span></div>
\end{htmlcode}
\choices{\texttt{10 + 20}}{\texttt{30}}{\texttt{30.0}}{\texttt{1020}}

\noindent\textbf{Câu 14.} Quan sát câu lệnh sau:
\begin{htmlcode}
<div style="background-color: #0f172a; color: #ffffff; padding: 10px 16px; margin: 6px 0; border-radius: 6px; border: 1.5px solid #3b82f6; font-family: monospace; font-size: 14px;"><span style="color: #38bdf8; font-weight: bold;">print</span><span style="color: #ffffff; font-weight: bold;">(</span><span style="color: #4ade80; font-weight: bold;">13</span> <span style="color: #f472b6; font-weight: bold;">+</span> <span style="color: #4ade80; font-weight: bold;">10</span> <span style="color: #f472b6; font-weight: bold;">*</span> <span style="color: #4ade80; font-weight: bold;">3</span> <span style="color: #f472b6; font-weight: bold;">/</span> <span style="color: #4ade80; font-weight: bold;">2</span> <span style="color: #f472b6; font-weight: bold;">-</span> <span style="color: #4ade80; font-weight: bold;">3</span> <span style="color: #f472b6; font-weight: bold;">*</span> <span style="color: #4ade80; font-weight: bold;">2</span><span style="color: #ffffff; font-weight: bold;">)</span></div>
\end{htmlcode}
Kết quả in ra màn hình của biểu thức trên là:
\choices{\texttt{22.0}}{\texttt{22}}{\texttt{28.5}}{\texttt{16.0}}

\noindent\textbf{Câu 15.} Cho câu lệnh:
\begin{htmlcode}
<div style="background-color: #0f172a; color: #ffffff; padding: 10px 16px; margin: 6px 0; border-radius: 6px; border: 1.5px solid #3b82f6; font-family: monospace; font-size: 14px;"><span style="color: #38bdf8; font-weight: bold;">print</span><span style="color: #ffffff; font-weight: bold;">(</span><span style="color: #4ade80; font-weight: bold;">9</span> <span style="color: #f472b6; font-weight: bold;">/</span> <span style="color: #4ade80; font-weight: bold;">3</span><span style="color: #ffffff; font-weight: bold;">)</span></div>
\end{htmlcode}
Kiểu dữ liệu của giá trị in ra màn hình là gì?
\choices{Kiểu số thực (\texttt{float})}{Kiểu số nguyên (\texttt{int})}{Kiểu xâu ký tự (\texttt{str})}{Kiểu lôgic (\texttt{bool})}

\noindent\textbf{Câu 16.} Trong Python, hàm nào sau đây được dùng để xuất (in) dữ liệu hoặc thông báo ra màn hình?
\choices{\texttt{print()}}{\texttt{input()}}{\texttt{output()}}{\texttt{write()}}

\noindent\textbf{Câu 17.} Python là ngôn ngữ lập trình thuộc loại nào sau đây?
\choices{Ngôn ngữ lập trình bậc cao}{Ngôn ngữ máy}{Hợp ngữ}{Ngôn ngữ nhị phân}

\noindent\textbf{Câu 18.} Người sáng lập và cho ra mắt ngôn ngữ lập trình Python lần đầu tiên vào năm 1991 là ai?
\choices{Guido van Rossum}{Tim Berners-Lee}{Bill Gates}{James Gosling}

\noindent\textbf{Câu 19.} Để máy tính có thể hiểu và thực thi mã nguồn viết bằng ngôn ngữ bậc cao Python, máy tính cần sử dụng thành phần nào?
\choices{Chương trình dịch (trình thông dịch)}{Trình duyệt web}{Thiết bị lưu trữ ngoài (USB)}{Bộ nhớ RAM}

\noindent\textbf{Câu 20.} Phần mở rộng mặc định của tệp lưu mã nguồn chương trình Python là gì?
\choices{\texttt{.py}}{\texttt{.python}}{\texttt{.pt}}{\texttt{.doc}}

\end{document}
